\chapter{Valores de referência}


Os valores de referência constituem parâmetros técnicos utilizados como base comparativa para avaliar o desempenho energético da unidade consumidora. A partir dessa comparação, é possível estimar o potencial de economia de energia com a adoção de medidas de eficiência energética e classificar o nível de desempenho atual da edificação frente a padrões consolidados.

A Tabela~\ref{tab:valores-referencia} apresenta, a título de comparação, valores típicos de consumo específico por tipo de edificação, úteis para análise do desempenho energético da unidade em estudo.


\begin{table}[H]
\centering
\footnotesize
\renewcommand{\arraystretch}{1.3}
\begin{tabular}{|>{\columncolor[HTML]{1a3a5c}\color{white}}p{6cm}|p{6cm}|}
\hline
\textbf{Tipo de edificação} & \textbf{Consumo específico típico} \\ \hline
<<tipoEdificacao>>
& <<consumoEspecificoTipico>> \\ \hline
\end{tabular}
\caption{Valores típicos de consumo específico por tipo de edificação pública}
\label{tab:valores-referencia}
\vspace{0.3em}
\begin{flushleft}
\footnotesize * Os valores são indicativos e podem variar conforme fatores climáticos, regime de operação, ocupação, tipo de equipamentos e práticas de manutenção.
\end{flushleft}
\end{table}

A comparação entre os indicadores reais da unidade e os valores de referência apresentados permite identificar desvios de desempenho, avaliando o grau de eficiência energética atual da edificação. Essa análise fornece uma base objetiva para a estimativa do potencial técnico de economia de energia, orientando a priorização das medidas propostas. Além disso, constitui um elemento fundamental para o dimensionamento das metas de economia a serem alcançadas com a implementação das ações, conforme estabelecido nos critérios metodológicos definidos no \textbf{Módulo 7 do Programa de Eficiência Energética da ANEEL (PROPEE)}~\cite{propee_aneel}, que tratam da quantificação de resultados e do cálculo da relação custo-benefício dos projetos.

\vspace{1cm}

A Tabela~\ref{tab:indicadores-desempenho} apresenta a avaliação dos índices de desempenho da unidade consumidora frente aos respectivos valores de \textit{benchmarking} adotados como referência.

\begin{table}[H]
\centering
\footnotesize
\renewcommand{\arraystretch}{1.3}
\begin{tabular}{|>{\columncolor[HTML]{1a3a5c}\color{white}}p{5.5cm}|p{3.5cm}|p{4.5cm}|}
\hline
\textbf{Índices} & \textbf{Média Ano} & \textbf{Benchmarking} \\ \hline
Fator de Carga na Ponta            & <<fatorCargaPonta>>
         & Maior que 0,80           \\ \hline
Fator de Carga Fora Ponta          & <<fatorCargaForaPonta>>
        & Maior que 0,60           \\ \hline
Fator de Potência                  & <<fatorPotencia>>
         & Maior que 0,92           \\ \hline
Custo Médio Unitário (R\$/kWh)     & R\$ <<custoMedioUnitario>>
     & Menor que 0,60           \\ \hline
Consumo por <<parametro>> (kWh/mês)        & <<consumoPessoa>>
 kWh/mês    & Menor que 7 kWh/mês      \\ \hline
\end{tabular}
\caption{Avaliação dos índices de desempenho da unidade consumidora}
\label{tab:indicadores-desempenho}
\end{table}

A análise conjunta dos indicadores permite identificar os seguintes pontos:

\vspace{1cm}

\noindent\textbf{Pontos Fortes:}
\begin{itemize}
    \item \textbf{Fator de Potência adequado} $\rightarrow$ sem risco de multa (Resolução ANEEL 1.000/2021).
    \item \textbf{Consumo por Área e por Funcionário abaixo dos \textit{benchmarks}} $\rightarrow$ boa gestão estrutural e operacional.
\end{itemize}

\noindent\textbf{Pontos Críticos:}
\begin{itemize}
    \item \textbf{Fator de Carga (ponta e fora ponta) muito abaixo do ideal} $\rightarrow$ indica operação concentrada em poucas horas e/ou equipamentos subutilizados.
    \item \textbf{Custo médio do kWh elevado (R\$ 0,73)} $\rightarrow$ possível inadequação tarifária e/ou de demanda contratada e/ou horário de ponta mal gerenciado.
    \item \textbf{Consumo por Aluno elevado} $\rightarrow$ desperdício associado a atividades diurnas (luz acesa à toa, ar-condicionado mal regulado, equipamentos em \textit{stand-by}).
\end{itemize}

De forma geral, a unidade consumidora (\textbf{<<nomeUc>>}) apresenta eficiência estrutural boa (consumo por área e por funcionário), mas sofre de \textbf{ineficiência operacional} (baixo fator de carga e alto custo por kWh e por aluno).

Existe, portanto, potencial técnico de economia que pode ser alcançado predominantemente com medidas de gerenciamento de cargas de climatização e iluminação, programação de horários de funcionamento e revisão do enquadramento tarifário.


% Exemplo de tabela de referência (caso queira inserir valores típicos):
% \begin{table}[H]
% \centering
% \renewcommand{\arraystretch}{1.3}
% \begin{tabular}{|>{\columncolor{azuluc}\color{white}}l|c|}
% \hline
% \textbf{Tipo de Edificação} & \textbf{Consumo Específico Típico (kWh/m²/mês)} \\ \hline
% Escolas públicas & 4,0 a 8,0 \\ \hline
% Unidades de saúde & 10,0 a 20,0 \\ \hline
% Prédios administrativos & 6,0 a 12,0 \\ \hline
% \end{tabular}
% \caption{Valores típicos de consumo específico por tipo de edificação}
% \label{tab:valores-referencia}
% \end{table}
